% !TeX root=../main.tex
% در این فایل، عنوان پایان‌نامه، مشخصات خود، متن تقدیمی‌، ستایش، سپاس‌گزاری و چکیده پایان‌نامه را به فارسی، وارد کنید.
% توجه داشته باشید که جدول حاوی مشخصات پروژه/پایان‌نامه/رساله و همچنین، مشخصات داخل آن، به طور خودکار، درج می‌شود.
%%%%%%%%%%%%%%%%%%%%%%%%%%%%%%%%%%%%
% دانشگاه خود را وارد کنید
\university{دانشگاه تهران}
% پردیس دانشگاهی خود را اگر نیاز است وارد کنید (مثال: فنی، علوم پایه، علوم انسانی و ...)
\college{پردیس دانشکده‌های فنی}
% دانشکده، آموزشکده و یا پژوهشکده  خود را وارد کنید
\faculty{دانشکدهٔ علوم مهندسی}
% گروه آموزشی خود را وارد کنید (در صورت نیاز)
\department{گروه کنترل}
% رشته تحصیلی خود را وارد کنید
\subject{مهندسی یرق}
% گرایش خود را وارد کنید
\field{کنترل}
% عنوان پایان‌نامه را وارد کنید
\title{طراحی و ساخت حسگر لامسه ی منعطف بر اساس ترکیب داده های حسگری برای تعامل و کار کردن با اشیاء}
% نام استاد(ان) راهنما را وارد کنید
\firstsupervisor{دکتر مهدی طالع ماسوله}
\firstsupervisorrank{دانشیار}
\secondsupervisor{دکتر احمد کلهر}
\secondsupervisorrank{دانشیار}
% نام استاد(دان) مشاور را وارد کنید. چنانچه استاد مشاور ندارید، دستورات پایین را غیرفعال کنید.
\firstadvisor{دکتر محمدرضا نیری}
\firstadvisorrank{استادیار}
%\secondadvisor{دکتر مشاور دوم}
% نام داوران داخلی و خارجی خود را وارد نمایید.
\internaljudge{دکتر شاهین جعفرآبادی آشتیانی}
\internaljudgerank{دانشیار}
\externaljudge{دکتر محمداعظم خسروی}
\externaljudgerank{دانشیار}
\externaljudgeuniversity{دانشگاه صنعتی امیرکبیر}
% نام نماینده کمیته تحصیلات تکمیلی در دانشکده \ گروه
\graduatedeputy{دکتر نماینده}
\graduatedeputyrank{دانشیار}
% نام دانشجو را وارد کنید
\name{محمدرضا}
% نام خانوادگی دانشجو را وارد کنید
\surname{شیخ عظیمی پور سردرود}
% شماره دانشجویی دانشجو را وارد کنید
\studentID{810101341}
% تاریخ پایان‌نامه را وارد کنید
\thesisdate{شهریور 1405}
% به صورت پیش‌فرض برای پایان‌نامه‌های کارشناسی تا دکترا به ترتیب از عبارات «پروژه»، «پایان‌نامه» و «رساله» استفاده می‌شود؛ اگر  نمی‌پسندید هر عنوانی را که مایلید در دستور زیر قرار داده و آنرا از حالت توضیح خارج کنید.
%\projectLabel{پایان‌نامه}

% به صورت پیش‌فرض برای عناوین مقاطع تحصیلی کارشناسی تا دکترا به ترتیب از عبارت «کارشناسی»، «کارشناسی ارشد» و «دکتری» استفاده می‌شود؛ اگر نمی‌پسندید هر عنوانی را که مایلید در دستور زیر قرار داده و آنرا از حالت توضیح خارج کنید.
%\degree{}
%%%%%%%%%%%%%%%%%%%%%%%%%%%%%%%%%%%%%%%%%%%%%%%%%%%%
%% پایان‌نامه خود را تقدیم کنید! %%
\dedication
{
{\Large تقدیم به:}\\
\begin{flushleft}{
	\huge
	خانواده‌ام\\
	\vspace{7mm}
	حامیان و یاوران همیشگی
}
\end{flushleft}
}
%% متن قدردانی %%
%% ترجیحا با توجه به ذوق و سلیقه خود متن قدردانی را تغییر دهید.
\acknowledgement{
با سپاس بیکران به درگاه خداوند متعال که توفیق و توانایی پیمودن این مسیر علمی را به من عطا فرمود. وظیفه خود می‌دانم از استاد بزرگوارم، جناب آقای دکتر مهدی طالع ماسوله، به عنوان استاد راهنمای اصلی این پایان‌نامه، صمیمانه‌ترین سپاس‌ها را ابراز نمایم. ایشان نه تنها با راهنمایی‌های علمی و نکات ارزشمندشان مسیر پژوهش را برایم روشن ساختند، بلکه با ایجاد فضایی صمیمی و پویا در محیط علمی، انگیزه و دلگرمی فراوانی در من برانگیختند. حمایت‌های مادی و معنوی ایشان، به‌ویژه پشتیبانی مالی از این پژوهش، سهم بسزایی در شکل‌گیری و به ثمر رسیدن این کار داشته است. بی‌تردید بخش مهمی از موفقیت این پایان‌نامه مرهون همراهی، بزرگواری و حمایت‌های بی‌دریغ ایشان است. همچنین از جناب آقای دکتر احمد کلهر، استاد راهنمای دوم، به خاطر راهنمایی‌های ارزشمند و حمایت‌های علمی و اخلاقیشان سپاسگزارم. حضور ایشان در کنار این مسیر، همواره موجب دلگرمی غنای علمی کار بوده است. همچنین قدردانی ویژه خود را تقدیم می‌دارم به جناب آقای دکتر محمدرضا نیری، استاد مشاور محترم، که با دانش فنی عمیق، نگاه دقیق و نکات ارزشمندشان، نقش بی‌بدیلی در پیشبرد علمی این پژوهش ایفا کردند. بدون راهنمایی‌ها و پیشنهادهای ایشان، دستیابی به بسیاری از نتایج این کار امکان‌پذیر نبود.
در ادامه، لازم می‌دانم از خانواده عزیزم، به‌ویژه پدر و مادر مهربانم و خواهر عزیزم، قدردانی کنم. صبر، محبت، فداکاری و حمایت بی‌دریغ آنان در تمام مراحل زندگی و تحصیل، بزرگ‌ترین پشتوانه من بوده است. با افتخار و از صمیم قلب، بوسه بر دستان پرمهر پدر و مادرم می‌زنم، چرا که هر موفقیتی که امروز به آن دست یافته‌ام، مرهون زحمات، محبت‌ها و دعای خیر خانواده‌ام است.
همچنین از دوستان و همکاران نازنین و گرانقدرم در آزمایشگاه تعامل انسان و ربات که با همدلی و همراهی، همواره مشوق و پشتیبان من در این مسیر بودند، صمیمانه سپاسگزارم. همواره عضو کوچکی از این جمع غنی و مترقی خواهم بود، چرا آشنایی با این عزیزان از بهترین اتفاقات حیات من است.
در پایان، این اثر را تقدیم می‌کنم به تمامی عزیزانی که در این راه دشوار و شیرین در کنارم بودند و به هر شکل یاری‌گر من شدند. امید است حاصل این تلاش، گامی هرچند کوچک در اعتلای دانش و خدمت به جامعه باشد.
}
%%%%%%%%%%%%%%%%%%%%%%%%%%%%%%%%%%%%
%چکیده پایان‌نامه را وارد کنید
\fa-abstract{
این پژوهش به طراحی، ساخت، کالیبراسیون و ارزیابی یک حسگر لامسه چندوجهی و نرم، با الهام از زیست و انسان، برای افزایش قابلیت‌های تعاملی ربات‌ها می‌پردازد. با الهام از ماهیت چندوجهی حس لامسه انسان که اطلاعات استاتیکی و دینامیکی را به صورت مجزا دریافت می‌کند. این پژوهش رویکردی نوین برای غلبه بر محدودیت‌های حسگرهای تک‌وجهی ارائه می‌دهد که ابتدا با طراحی و ساخت یک نمونه اولیه تک‌‌محوره مبتنی بر حسگر فشار بارومتریک آغاز شد که ضمن اثبات کارایی این مکانیزم در اندازه‌گیری دقیق نیروی عمودی، خواص خطی‌بودن و پسماند نیز ارزیابی شد. از آنجایی که این نمونه قابلیت اندازه‌گیری نیرو های مماسی و ماهیت چندوجهی را نداشته، یک مجموعه‌ی چندوجهی یکپارچه توسعه یافت که شامل یک حسگر بارومتریک، دو حسگر اثر هال به همراه یک آهنربای داخلی برای تفکیک نیروهای عمودی و مماسی، و یک میکروفون ممز به عنوان حسگر پیزوالکتریک برای ثبت ارتعاشات فرکانس بالا ناشی از لغزش و بافت سطح بود. به منظور مدیریت و همگام‌سازی حجم بالای داده‌های ناهمگون، یک برد جمع‌آوری داده اختصاصی با قابلیت‌های نوآورانه نظیر پروتکل ارتباطی ترکیبی USB و بهره‌گیری پیشرفته از دسترسی مستقیم به حافظه نیز طراحی شد. در ادامه، به دلیل ماهیت غیرخطی و پیچیده‌ی حسگر چند‌وجهی، فرآیند کالیبراسیون با استفاده از مدل‌های یادگیری عمیق انجام پذیرفت؛ یک شبکه عصبی تماماً متصل برای مدل‌سازی رفتار استاتیک و استخراج دقیق بردارهای نیرو، و یک شبکه حافظه‌دار برای شناسایی ماهیت دینامیکی و وابسته به زمان مجموعه‌ حسگر آموزش داده شد. دقت مجموعه‌ی ارائه شده برای نیروی عمودی $99.75\%$ و برای نیرو مماسی $98.1\%$ به دست آمد. در نهایت، کارایی کل سامانه در یک پیاده‌سازی عملی به اثبات رسید؛ حسگر بر روی یک چنگک رباتیک نصب شد و در یک حلقه کنترلی، قابلیت کنترل دقیق نیروی گرفتن همزمان با تشخیص پدیده های دینامیکی بر اساس داده‌های ارتعاشی با موفقیت به نمایش درآمد.در مجموع، این دستاوردها یک چارچوب سخت‌افزاری و نرم‌افزاری جامع برای ارتقاء درک لامسه رباتیک و دستیابی به تعاملات هوشمند و ایمن با اشیاء ارائه داده است.

}
\keywords{حسگر چندوجهی، زیست‌الهام، حسگر لامسه، کنترل نیرو، تشخیص لغزش}

% انتهای وارد کردن فیلد‌ها
%%%%%%%%%%%%%%%%%%%%%%%%%%%%%%%%%%%%%%%%%%%%%%%%%%%%%%
